\section*{الفصل \ref{ch:b3:complex-mechanisms} --- آليات تفاعلات المعقدات}

\begin{solution}{exo:b3:complex-mechanisms:1}
خاملان: $\ce{[Cr(H2O)6]^{3+}}$ ($d^3$) و$\ce{[Co(NH3)6]^{3+}}$ ($d^6$ منخفض السبين)، بطاقتي تنشيط
كبيرتين لمجال الربائط. سريعة الاستبدال: $\ce{[Cr(H2O)6]^{2+}}$ ($d^4$ عالي السبين) و$\ce{[Cu(H2O)6]^{2+}}$ ($d^9$)،
المشوَّهان بمفعول يان--تيلر بروابط محورية طويلة وضعيفة؛ و$\ce{[Zn(H2O)6]^{2+}}$ ($d^{10}$)، بلا حاجز
لمجال الربائط.
\end{solution}

\begin{solution}{exo:b3:complex-mechanisms:2}
حين $k_2[\mathrm Y] \gg k_{-1}[\mathrm X]$ يكون المقام $k_2[\mathrm Y]$ و$v = k_1[\mathrm{ML_5X}]$؛ وحين $k_{-1}[\mathrm X] \gg k_2[\mathrm Y]$ يكون
$k_{-1}[\mathrm X]$ وتنتج الصيغة الثانية (توازن مسبق تثبّطه الربيطة
المغادرة).
\end{solution}

\begin{solution}{exo:b3:complex-mechanisms:3}
D: $\Delta V^{\ddagger} > 0$ و$\Delta^{\ddagger}S^\circ > 0$ (تغادر ربيطة). A: كلاهما سالب (ترتبط ربيطة في
الحالة الانتقالية).
\end{solution}

\begin{solution}{exo:b3:complex-mechanisms:4}
\babelsublr{$cis$-$\ce{[PtCl2(NH3)2]}$} انطلاقًا من $\ce{[PtCl4]^{2-}}$؛ و\babelsublr{$trans$-$\ce{[PtCl2(NH3)2]}$} انطلاقًا من $\ce{[Pt(NH3)4]^{2+}}$ (\omterm{def:b3:complex-mechanisms:trans-effect}{فمفعول
ترانس} للأيون \ce{Cl-} يفوق مفعول \ce{NH3}).
\end{solution}

\begin{solution}{exo:b3:complex-mechanisms:5}
$k_{\mathrm{obs}} = 2.0 \times \num{3.0e4} \times 0.10/1.2 = \qty{5.0e3}{s^{-1}}$؛ وعند \qty{1.0}{mol/L}، $\num{6.0e4}/3.0 = \qty{2.0e4}{s^{-1}}$؛ والحد
$k_i = \qty{3.0e4}{s^{-1}}$.
\end{solution}

\begin{solution}{exo:b3:complex-mechanisms:6}
$\Delta V^{\ddagger} = -RT\ln2/\Delta p = -8.314 \times 298 \times 0.693/\num{99.9e6} = \qty{-1.7e-5}{m^3.mol^{-1}} = \qty{-17}{cm^3.mol^{-1}}$:
\omterm{def:b3:complex-mechanisms:mechanisms}{فالآلية ترابطية}.
\end{solution}

\begin{solution}{exo:b3:complex-mechanisms:7}
الفوسفين، المانح $\sigma$ القوي والمستقبل $\pi$، يُضعف الرابطة المقابلة له (\omterm{def:b3:complex-mechanisms:trans-effect}{تأثير
ترانس}، رابطة أطول) ويجعلها سريعة الاستبدال (\omterm{def:b3:complex-mechanisms:trans-effect}{مفعول ترانس}): فذلك الكلوريد هو الذي
يُستبدل أولًا.
\end{solution}

\begin{solution}{exo:b3:complex-mechanisms:8}
تنتهي الربيطة الجسرية منقولةً إلى الفلز المؤكسَد؛ وتتوقف السرعة
بشدة على طبيعة الجسر المحتمل؛ ويُكشف وسيط ثنائي النواة؛ ولا يمكن أن
تتجاوز السرعة سرعة الاستبدال اللازم لتكوين
الجسر على الشريك \omterm{def:b3:complex-mechanisms:labile}{السريع الاستبدال}.
\end{solution}

\begin{solution}{exo:b3:complex-mechanisms:9}
$(\lambda + \Delta G^\circ)^2/4\lambda$ مع $4\lambda = 320$: 20.0 و11.3 و0 و\qty{5.0}{kJ/mol} (الأخيرة في \omterm{def:b3:complex-mechanisms:inverted}{المنطقة
المقلوبة}).
\end{solution}

\begin{solution}{exo:b3:complex-mechanisms:10}
$k_{12} = \sqrt{4.0 \times \num{2.0e5} \times \num{1.0e4}} = \qty{8.9e4}{L.mol^{-1}.s^{-1}}$.
\end{solution}

\begin{solution}{exo:b3:complex-mechanisms:11}
$d^5$ عالي السبين: 0؛ $d^7$: $1.33\,Dq$؛ $d^8$: $2.0\,Dq$. وسرعات تبادل الماء المتوقعة: \ce{Mn^{2+}} $>$ \ce{Co^{2+}} $>$
\ce{Ni^{2+}}، وهو الترتيب المرصود.
\end{solution}

\begin{solution}{exo:b3:complex-mechanisms:12}
في التفاعلات الشديدة الإطلاق لطاقة جيبس تتجاوز السرعة الذاتية سرعة التقاء
الأيونات: فيبقى الثابت الملاحظ عند حد الانتشار ويختفي
التناقص. ومع تثبيت المانح والمستقبل على مسافة محددة في جزيء واحد، لا
توجد خطوة انتشار، وأظهرت سلسلة من المستقبلات ذات القوة الدافعة المتزايدة
أن السرعة ترتفع، وتمر بقيمة عظمى ثم تنخفض.
\end{solution}

\begin{solution}{pb:b3:complex-mechanisms:1}
\textbf{1.} $5d^8$.
\textbf{2.} المجال القوي لأيون $5d$ يجعل المدار $d_{x^2-y^2}$، المتجه نحو أربع
ربائط، مرتفعًا جدًا: فتزدوج الإلكترونات الثمانية في المدارات الأربعة الدنيا.
\textbf{3.} للمعقد 16 إلكترون تكافؤ وموقع محوري مفتوح: فترتبط الربيطة
الداخلة أولًا، عبر حالة انتقالية خماسية التناسق.
\textbf{4.} $k_{\mathrm{obs}} = k_1 + k_2[\mathrm Y]$: طريق المذيب (هجوم المذيب، ثم استبداله السريع بالربيطة Y)
وطريق مباشر.
\textbf{5.} يستغرق الاستبدال ساعات (الجزء الثالث): فهو خامل على مقياس الدقائق، لا على مقياس
اليوم.
\textbf{6.} السيسبلاتين، لأن جزيء \ce{NH3} الثاني يحل محل كلوريد مقابل لكلوريد
(\omterm{def:b3:complex-mechanisms:trans-effect}{مفعول ترانس} $\ce{Cl-} > \ce{NH3}$)؛ وتتكوّن نواتج جانبية أيضًا.
\textbf{7.} لليوديد \omterm{def:b3:complex-mechanisms:trans-effect}{مفعول ترانس} أكبر من الكلوريد، وهذا يجعل التوجيه
أنظف وأسرع.
\textbf{8.} في $\ce{[PtI3(NH3)]-}$ يصير اليوديد المقابل لليوديد \omterm{def:b3:complex-mechanisms:labile}{سريع الاستبدال}، لا المقابل
للجزيء \ce{NH3}: فيدخل جزيء الأمونيا الثاني في موضع cis من الأول.
\textbf{9.} يترسب يوديد الفضة ويبقى \babelsublr{$cis$-$\ce{[Pt(H2O)2(NH3)2]^{2+}}$} في المحلول.
\textbf{10.} يحل الكلوريد محل جزيئي الماء في موضعيهما نفسيهما؛ ولا
تُمس ربائط الأمين.
\textbf{11.} انطلاقًا من $\ce{[Pt(NH3)4]^{2+}}$ ومن أيوني كلوريد.
\textbf{12.} موضعاه النشطان متقابلان: فلا يستطيع أن يرتبط بقاعدتين
متجاورتين من سلسلة الحمض النووي نفسها، وهي الإصابة التي يعمل بها
السيسبلاتين.
\textbf{13.} $\ln2/\num{9.0e-5} = \qty{7.7e3}{s} = \qty{2.1}{h}$.
\textbf{14.} $1 - \eu^{-\num{9.0e-5} \times 7200} = 0.48$.
\textbf{15.} الماء مجموعة مغادرة أفضل بكثير من الكلوريد: فيُستبدل المعقد المائي
بسرعة بذرة نيتروجين لقاعدة من قواعد الحمض النووي.
\textbf{16.} ترابطي، كسائر استبدالات البلاتين(II): \omterm{def:b3:complex-mechanisms:activation-volume}{حجم تنشيط}
سالب.
\textbf{17.} يدفع تركيز الكلوريد المرتفع التوازن عائدًا نحو الصورة
الثنائية الكلور، فيبقى الدواء سليمًا حتى يبلغ الخلايا.
\textbf{18.} $f = K/(K + [\ce{Cl-}])$.
\textbf{19.} $\num{3.0e-3}/0.103 = 0.029$.
\textbf{20.} $\num{3.0e-3}/0.0070 = 0.43$.
\textbf{21.} 15.
\textbf{22.} لا يوجد جزء معتبر في صورة المعقد المائي النشط إلا داخل
الخلايا.
\textbf{23.} الربائط الكبريتية، مثل الغلوتاثيون والسيستئين والميثيونين في
البروتينات، التي ترتبط بقوة بالبلاتين اللين.
\textbf{24.} ترتبط الأمونيا بالبلاتين بقوة وتقع مقابل الكلوريد، وهو ربيطة
\omterm{def:b3:complex-mechanisms:trans-effect}{مفعول ترانس} فيها صغير: فلا تصير روابط الأمين سريعة الاستبدال.
\textbf{25.} \textbf{يوجد من الدواء في الصورة المائية النشطة داخل الخلية نحو 15 مرة أكثر
مما في البلازما.}
\end{solution}
